


\section{Related Work}

% \noindent\textbf{Contrastive Vision Language Models.}  
% Early multimodal representation learning approaches mapped latent representations of text and images into a shared space by aligning disjoint encoders through contrastive learning~\citep{radford2021learningtransferablevisualmodels,zhai2023sigmoidlosslanguageimage}. These models laid the foundation for visual embedding tasks such as classification and image caption retrieval. However, they are not optimized for applications that require deeper visual document understanding, as they lack token-level cross-modality alignment objectives and face architectural limitations~\citep{faysse2025colpaliefficientdocumentretrieval,jiang2024e5vuniversalembeddingsmultimodal,jiang2025vlm2vectrainingvisionlanguagemodels}.

% \noindent\textbf{Generative Vision Language Models.}  
% Building on advances of large language models (LLMs), recent work has combined transformer-based text backbones with contrastively pretrained ViT encoders to create large generative vision–language models (VLMs)~\citep{li2022blipbootstrappinglanguageimagepretraining,alayrac2022flamingovisuallanguagemodel,bai_qwen-vl_2023,deitke2024molmopixmoopenweights,wang2024qwen2vlenhancingvisionlanguagemodels,yang2025qwen251mtechnicalreport}. These models leverage LLM attention to enrich image–text understanding and instruction-following, advancing applications such as object detection, OCR, and image captioning, and establishing themselves as core to the current multimodal foundation model landscape~\citep{alayrac2022flamingovisuallanguagemodel,beyer2024paligemmaversatile3bvlm,nassar2025smoldoclingultracompactvisionlanguagemodel}. Nevertheless, generative VLMs are often prone to visual hallucinations~\citep{tong2024eyeswideshutexploring} and are optimized for next-token prediction rather than general-purpose embeddings~\citep{jiang2025vlm2vectrainingvisionlanguagemodels}, making them an imperfect fit for retrieval-oriented applications.

\noindent\textbf{Repurposing VLMs for Representation Learning.}  
Motivated by the zero-shot performances of generative VLMs~\citep{alayrac2022flamingovisuallanguagemodel, lucas_beyer_paligemma_2024, bai_qwen-vl_2023}, recent studies have explored repurposing these for multimodal embedding tasks~\citep{ma2024unifyingmultimodalretrievaldocument,faysse2025colpaliefficientdocumentretrieval,jiang2025vlm2vectrainingvisionlanguagemodels, zhang2025gmeimprovinguniversalmultimodal}. As backbone generative models improved, retriever performance improved as well showcasing the central impact of language model pretraining and modality alignment \citep{xu2025llamanemoretrievercolembedtopperforming, nussbaum2025nomicembedtrainingreproducible}. 
% For example, VLM2Vec~\citep{jiang2025vlm2vectrainingvisionlanguagemodels} is based on Phi-3.5-V~\citep{abdin_phi-3_2024}, while MoCa on Qwen-2.5-VL~\citep{bai2025qwen25vltechnicalreport}. 
These model remain inherently constrained by their causal attention mechanisms which has been shown in text settings to limits represational efficiency ~\citep{gisserotboukhlef2025pretrainencodersmaskedlanguage,weller2025seqvsseqopen}. Recent work attempts to address this issue by modifying VLM attention during continual pretraining \citep{chen2025mocamodalityawarecontinualpretraining} or contrastive tuning ~\citep{jiang2025vlm2vectrainingvisionlanguagemodels, xu2025llamanemoretrievercolembedtopperforming}, but no recent work attempts to align natively bidirectional language encoder models with vision encoders. The recent release of long sequence text encoders \citep{modernbert, boizard2025eurobertscalingmultilingualencoders} makes this possible.
% These approaches achieve notable gains on retrieval benchmarks, largely due to the contextualization power of LLM backbones, and on the inherent capabilities of VLMs to address complex visual tasks, including optical character recognition (OCR) and visual question answering (VQA)~\citep{wang2024qwen2vlenhancingvisionlanguagemodels,internvl2-5,nassar2025smoldoclingultracompactvisionlanguagemodel}. 
% Yet, they remain inherently constrained: causal attention mechanisms limit their efficiency~\citep{gisserotboukhlef2025pretrainencodersmaskedlanguage,chen2025mocamodalityawarecontinualpretraining}, and their large scale often makes them impractical for industrial deployment~\citep{marafioti2025smolvlmredefiningsmallefficient}. Recent work attempts to address this issue by modifying VLM attention through continual pretraining~\cite{chen2025mocamodalityawarecontinualpretraining}, but does not assess the benefits of training a vision retriever from scratch with native bidirectional attention. 

\noindent\textbf{Late Interaction in Visual Document Retrieval}  To further boost performance, visual document retrievers leverage the late interaction mechanism~\citep{khattab_colbert_2020} which matches multiple query embeddings with multiple document embeddings through the MaxSim operation ~\citep{faysse2025colpaliefficientdocumentretrieval,günther2025jinaembeddingsv4universalembeddingsmultimodal, xu2025llamanemoretrievercolembedtopperforming}. This enables more granular interactions between image and query tokens, at the cost of additional storage and a slight compute overhead during the matching operation. Efficiency gains have come from improving the storage costs through quantization \citep{vespaScalingColPali}, token pruning \citep{faysse_croissantllm_2024} and more recently the use of Matrioshka losses to compact multi-token representations \citep{xiao2025metaembedscalingmultimodalretrieval}. Ultimately, the performance bottleneck when running visual retrieval inference with such models now resides mostly in the necessity to rely on costly GPU hardware to encode queries, which sets apart text from vision retrieval. This paper fills this gap by using encoders that run on CPU, of parameter sizes comparable to commonly used local text embedding models \citep{chen_bge_2024, enevoldsen2025mmtebmassivemultilingualtext}.

% trop specifgiaue, essaie toujours d englober plusoeurs papiers. ici ca peut etre genre token pooling poiur colbert (for efficiency) avec les matrioshka

% je t'avoue si t'as les papier en tete ca ira plus vite ahah surtt ce deuxieme paragraphe - ok je fais pas de soucis,m poli un oeu le reste (contrib 1) ok top cimer

%This tension between performance and efficiency motivates our investigation into whether large generative VLMs should remain the de facto choice for multimodal embedding.

%Faysse et al. specifically show that the OCR capabilities of VLMs, combined with the language understanding capabilities of the backbone language model enable creating high quality embeddings from document images that contain the semantic information contained in the document text. Many efforts since have optimized generative VLMs for embedding creation (VLM2Vec, VLM2Vec2, NemoRetriever, GTE).
